% ===================================================================
%  ТАБЛИЦЯ № 10 — Море. — Порт.
%  Traduction 2026 d'apres texto/io, controlee sur texto/fr, texto/en
%  et texto/es. Consigne : voir texto/uk/10-tablycia-01.tex.
%
%  Le francais decale sa numerotation d'un cran a partir de l'alinea 8.
%  C'est une coquille de l'imprimeur francais, conservee dans la
%  transcription ; l'ukrainien suit l'ido, qui compte juste.
%
%  LE LEXIQUE DE LA MER EST LE MOINS UKRAINIEN DE TOUS, et c'est une
%  chose a dire nettement : l'ukrainien n'a pas de tradition maritime
%  ancienne, et son vocabulaire de marine est venu tard, par les memes
%  chemins que celui du russe — le neerlandais et l'anglais, par
%  l'amirauté. « щогла », « палуба », « трап », « буксир » sont donc
%  les memes mots des deux cotes, et il n'y a pas de faute a les
%  ecrire ainsi. Ce qui reste ukrainien, ce sont les mots de la terre :
%  le brise-lames est un « хвилеріз », la douane une « митниця », le
%  camion une « вантажівка ».
% ===================================================================

%%K t10-apar-1 apar
\VUsaut{4.0mm}
\VUtitre{15.0pt}{60}{ТАБЛИЦЯ № 10}
\VUsaut{3.0mm}
\VUpk{11.4pt}{40}{Море. --- Порт.}
\VUsaut{3.0mm}

%%K t10-01-1 p
1. --- Улітку, поки одні шукають спокою і прохолоди в горах, інші воліють їхати
на узбережжя. Так ми й зробили торік, бо дуже любимо \VUgras{океан}
\textsuperscript{(1)}.

%%K t10-02-1 p
2. --- Що до мене, мені дає велику втіху дивитися на цю безмірну водяну
гладінь, простерту до \VUgras{обрію} \textsuperscript{(2)}, і дивитися, як
величезні \VUgras{пароплави} \textsuperscript{(3)} ідуть у довгий \VUgras{рейс} з
подорожніми, що прямують до Америки.

%%K t10-03-1 p
3. --- Тому батько, який знає мої смаки, завжди вибирає на канікули дуже тихий
берег біля одного з наших великих \VUgras{військових портів}.

%%K t10-03-2 p
Останнього нашого приїзду \VUgras{ескадра} прийшла і стала на якір за величезним
\VUgras{хвилерізом} \textsuperscript{(4)}, який замикає і боронить \VUgras{військову
гавань} \textsuperscript{(5)}, і я міг досхочу намилуватися різними
\VUgras{військовими суднами}: від маленького \VUgras{підводного човна}
\textsuperscript{(10)}, який пірнає і зникає під водою, до величезного
\VUgras{броненосця} \textsuperscript{(9)}, якому досі нічого було боятися, крім
нападів \VUgras{міноносця} \textsuperscript{(8)}.

%%K t10-tit-2 sub
\VUsaut{3.0mm}
\VUpk{10.2pt}{40}{Малі судна.}
\VUsaut{3.0mm}

%%K t10-04-1 p
4. --- Цей останній уже вмів \VUgras{спритно} знаходити слабке місце в товстій
металевій броні і ставити туди небезпечну \VUgras{торпеду}, вибух якої топить
судно; але його все-таки слід було боятися менше, ніж підводного човна, бо
винищувачі легко пускаються за ним у погоню.

%%K t10-05-1 p
5. --- Я міг бачити і швидкі броненосні \VUgras{крейсери} \textsuperscript{(6)},
майже такі самі невразливі, як справжній броненосець, кілька
\VUgras{транспортів} \textsuperscript{(12)}, три чи чотири \VUgras{канонерки}
\textsuperscript{(7)} і, нарешті, \VUgras{фрегат} \textsuperscript{(11)}.
\VUgras{Найцікавіший вигляд цього флоту, що знімається з якоря.}

%%K t10-06-1 p
6. --- Яка різниця в будові між цими громадами сталі і зграбними
\VUgras{тримачтовиками} чи граційними \VUgras{вітрильниками}
\textsuperscript{(13)}, які ще в ходу в торговельному флоті і які я бачив, як
вони входили в порт під усіма вітрилами, гуляючи по \VUgras{молу}
\textsuperscript{(14)}. Щоправда, вони багато втрачали, коли \VUgras{вітер} був
проти них. Доводилося прибирати всі вітрила, щоб увійти в док, і потрібен був
\VUgras{буксир} \textsuperscript{(16)}, щоб узяти їх і підвести до набережної, як
прості \VUgras{баржі} \textsuperscript{(17)}.

%%K t10-tit-3 sub
\VUsaut{3.0mm}
\VUpk{10.2pt}{40}{Торговельний порт.}
\VUsaut{3.0mm}

%%K t10-07-1 p
7. --- Порт, про який я кажу, тим і цікавий, що в ньому є не тільки військова
гавань, створена штучно велетенськими роботами, але й чудесний
\VUgras{торговельний порт} у \VUgras{гирлі} великої річки.

%%K t10-08-1 p
8. --- Смуга землі, що розділяє обидва доки, укріплена і боронена рядом
\VUgras{батарей} і \VUgras{бастіонів}, висунутих у море. Щоб гуляти по
\VUgras{валах} \textsuperscript{(15)}, потрібен окремий дозвіл. Я легко дістав
його через дядька, який служить інженером на \VUgras{верфях}
\textsuperscript{(18)}, і ходив дивитися, як під величезними повітками будують
судна.

%%K t10-09-1 p
9. --- Я бачив навіть спуск броненосця і пам'ятаю те хвилювання, яке огорнуло
весь натовп, коли величезна громада, полишена сама на себе, поповзла зі свого
ложа і зринула на воді річки.

%%K t10-10-1 p
10. --- Щоб потрапити на верфі, переходиш \VUgras{плац} \textsuperscript{(19)},
де щодня вправляються війська залоги, і йдеш залізним \VUgras{містком}
\textsuperscript{(20)}, що з'єднує плац з \VUgras{арсеналом}
\textsuperscript{(21)}.

%%K t10-tit-4 sub
\VUsaut{3.0mm}
\VUpk{10.2pt}{40}{Арсенал і доки.}
\VUsaut{3.0mm}

%%K t10-11-1 p
11. --- З дядьком я оглянув цей великий заклад, повний \VUgras{гармат}
\textsuperscript{(22)}, \VUgras{ядер} \textsuperscript{(23)} і
\VUgras{снарядів}. Бачив я і \VUgras{котельню} \textsuperscript{(24)}, де роблять
\VUgras{котли} \textsuperscript{(25)} для пароплавів.

%%K t10-12-1 p
12. --- Зовсім поруч --- \VUgras{доки} \textsuperscript{(26)}, сухі доки, і дуже
примітні \VUgras{плавучі доки} \textsuperscript{(27)}, у яких я міг зблизька
роздивитися на судні, що лагодиться, \VUgras{ґвинт} \textsuperscript{(28)},
\VUgras{кіль} \textsuperscript{(29)} і \VUgras{стерно} \textsuperscript{(30)}
корабля.

%%K t10-13-1 p
13. --- Торговельний порт вартий вивчення не менше за військовий, і я провів
чимало годин, тиняючись по набережних, де стоять \VUgras{колісні пароплави}
\textsuperscript{(31)}, які ходять угору річкою, і дивлячись, як працює
\VUgras{шеровий кран} \textsuperscript{(32)}, підіймаючи величезні тягарі, наче
пір'їнки.

%%K t10-tit-5 sub
\VUsaut{4.0mm}
\VUpk{10.2pt}{40}{Лоцман.}
\VUsaut{3.0mm}

%%K t10-14-1 p
14. --- Я милувався \VUgras{трансатлантичними пароплавами}
\textsuperscript{(34)}, що покидали рейд з розмаяними \VUgras{прапорами}
\textsuperscript{(35)}, яких вів умілий \VUgras{лоцман} \textsuperscript{(36)}, що
дивно вмів триматися \VUgras{фарватеру}, розміченого \VUgras{буями}
\textsuperscript{(40)} і \VUgras{віхами}, обминаючи \VUgras{ліхтери}
\textsuperscript{(41)}, якими так важко правити під єдиним чотирикутним
вітрилом. Я розрізняв на \VUgras{носі} \textsuperscript{(37)} \VUgras{каюти}
\textsuperscript{(39)} другого класу, а на \VUgras{кормі} \textsuperscript{(38)}
--- салони для пасажирів першого.

%%K t10-15-1 p
15. --- Інколи я дивився й на вантаження \VUgras{вантажного судна}
\textsuperscript{(42)}, у яке скидали товари зі \VUgras{складу}
\textsuperscript{(43)} і з \VUgras{портової станції} \textsuperscript{(44)}, що
їх підвозили численні потяги \VUgras{набережної колії} \textsuperscript{(45)}.

%%K t10-tit-6 sub
\VUsaut{3.0mm}
\VUpk{10.2pt}{40}{Катання на човні.}
\VUsaut{3.0mm}

%%K t10-16-1 p
16. --- Я часто катався на човні в \VUgras{закритому доку}
\textsuperscript{(53)}, де ховаються маленькі \VUgras{човни}
\textsuperscript{(46)}, і для мене була радість тримати \VUgras{весло}
\textsuperscript{(47)}. Я підіймався на \VUgras{палубу} \textsuperscript{(48)}
\VUgras{вітрильного човна} \textsuperscript{(49)}, вилазив по \VUgras{снастях}
\textsuperscript{(52)} на \VUgras{щоглу} \textsuperscript{(51)} до \VUgras{рей}
\textsuperscript{(50)}, а потім провадив прогулянку в \VUgras{ялі}
\textsuperscript{(54)} серед важких \VUgras{рибальських суден}
\textsuperscript{(55)}, де сохнуть розвішані \VUgras{сіті} \textsuperscript{(56)}.

%%K t10-17-1 p
17. --- На \VUgras{набережній} \textsuperscript{(58)} передо мною проходили
найрізноманітніші \VUgras{товари} \textsuperscript{(59)}, спаковані в
\VUgras{скрині} \textsuperscript{(60)} або в \VUgras{паки} \textsuperscript{(61)}.
Вони складали \VUgras{вантаж} \textsuperscript{(63)} барж, і потужні
\VUgras{крани} \textsuperscript{(62)} виймали їх і ставили на
\VUgras{вантажівки} \textsuperscript{(64)}, поки \VUgras{перевізники} оголошували
їх і платили мито в \VUgras{митниці} \textsuperscript{(65)}.

%%K t10-18-1 p
18. --- Нарешті, дійшовши до \VUgras{розвідного мосту} \textsuperscript{(66)} або
до \VUgras{шлюзу} \textsuperscript{(69)}, минувши \VUgras{землечерпалку}
\textsuperscript{(68)}, яка чистить \VUgras{канал} \textsuperscript{(67)}, я
приставав до молу біля \VUgras{сигнальної щогли} \textsuperscript{(70)}.

%%K t10-19-1 p
19. --- Саме з другого боку цього молу пристають \VUgras{катери}
\textsuperscript{(71)}, що привозять на берег офіцерів ескадри. Гарно дивитися,
як матроси разом підіймають весла, а шлюпка, несена \VUgras{хвилею}
\textsuperscript{(72)}, легко підходить до трапа. Звідси найкраще видно
\VUgras{приплив} і рух \VUgras{відпливу} та \VUgras{прибою} на \VUgras{пляжі}
\textsuperscript{(73)}.

%%K t10-tit-7 sub
\VUsaut{3.0mm}
\VUpk{10.2pt}{40}{Офіцерське зібрання.}
\VUsaut{3.0mm}

%%K t10-20-1 p
20. --- Щоб вернутися додому, я сідав в \VUgras{електричний трамвай}
\textsuperscript{(74)}, \VUgras{зупинка} \textsuperscript{(75)} якого коло
\VUgras{стовпа} перед офіцерським зібранням.

%%K t10-20-2 p
З \VUgras{балкона} \textsuperscript{(76)} зібрання, оздобленого \VUgras{якорями}
\textsuperscript{(77)}, гарматами і ядрами, я часто бачив блискуче товариство
морських офіцерів: \VUgras{лейтенантів} \textsuperscript{(78)} зі шпагами,
\VUgras{капітанів артилерії} \textsuperscript{(79)} і \VUgras{піхоти}
\textsuperscript{(80)} з \VUgras{шаблями}. Позаду них стояв \VUgras{матрос}
\textsuperscript{(81)}, який подавав їм \VUgras{підзорну трубу}, коли вони хотіли
роздивитися, що діється вдалині.

%%K t10-21-1 p
21. --- Бачив я і молодих \VUgras{мічманів} \textsuperscript{(82)}, які
діставали накази від свого \VUgras{командира} \textsuperscript{(83)} або від
\VUgras{адмірала} \textsuperscript{(84)}, поки \VUgras{квартирмейстер}
\textsuperscript{(85)} чекав, щоб відвезти їх на судно.

%%K t10-22-1 p
22. --- З цього балкона краєвид чудовий: він панує над усім пляжем з його
\VUgras{наметами} \textsuperscript{(86)} і \VUgras{купальними кабінами}
\textsuperscript{(87)}, і в годину купання видно, як \VUgras{купальники}
\textsuperscript{(88)} весело пустують перед \VUgras{казино}
\textsuperscript{(89)}.

%%K t10-23-1 p
23. --- У кінці пляжу берег утворює маленький скелястий \VUgras{півострів}
\textsuperscript{(91)}, об який розбиваються \VUgras{буруни}
\textsuperscript{(92)} і на якому збудовано \VUgras{маяк} \textsuperscript{(93)},
що вночі вказує морякам шлях. Цей півострів сполучений із землею лише дуже
вузьким \VUgras{перешийком} \textsuperscript{(90)}.

%%K t10-tit-8 sub
\VUsaut{3.0mm}
\VUpk{10.2pt}{40}{Загальний вигляд берега.}
\VUsaut{3.0mm}

%%K t10-24-1 p
24. --- \VUgras{Стрімчак} \textsuperscript{(94)} тягнеться далі, порізаний морем,
яке виточує в ньому химерну низку \VUgras{бухт} \textsuperscript{(95)} і
\VUgras{мисів}, роблячи його вельми мальовничим.

%%K t10-24-2 p
На \VUgras{плоскогір'ї} \textsuperscript{(96)}, що панує над \VUgras{протокою}
\textsuperscript{(98)}, стоять дуже важливий \VUgras{форт} \textsuperscript{(97)}
і кілька вілл.

%%K t10-25-1 p
25. --- Ця протока відділяє від суходолу дуже небезпечний \VUgras{острів}
\textsuperscript{(99)}, оточений \VUgras{скелями} \textsuperscript{(100)} і
\VUgras{рифами}, де інколи гинуть човни і навіть великі судна. Хоч яка прудка
поміч і хоч як швидко приходять \VUgras{рятувальні шлюпки}, ці \VUgras{корабельні
аварії} жахливі, і в рівноденні \VUgras{бурі} кілька суден загинуло з усією
залогою і вантажем.

%%K t10-25-2 p
Попри таке сусідство, \VUgras{порт} цей моряки відвідують охоче.
\VUsaut{6.0mm}
\VUfilet{22mm}
